% Chapter 5

\chapter{Conclusion and Next Steps}
\label{chap:conclusion}

This chapter concludes the \ac{PREPD} phase. The preparation work presented in this document establishes the foundation for thesis execution. No results, findings, or experimental outcomes are presented, as empirical work will be conducted during the subsequent dissertation phase.

\section{Synthesis of Preparation Work}

Chapter~\ref{chap:introduction} framed the research problem: observability in containerized \ac{CICD} platforms deployed on heterogeneous, on-premise infrastructure, historically constrained by limited operational access to the container runtime and still challenged by fragmentation, mixed Linux and Windows execution, and the absence of an integrated telemetry plane. The motivating gap is not only ``lack of root,'' but the distance between textbook cloud-native monitoring assumptions and enterprise CI/CD realities. The central research questions examine how an integrated observability framework can be designed and evaluated to improve reliability and reduce incident response times, including explicit trade-offs between instrumentation approaches. The introduction also documents an important evolution in the target organization: operations staff now have \texttt{docker} group access on Linux hosts, which relaxes several earlier practical constraints while leaving heterogeneity, change-management boundaries, and integration complexity as first-class concerns.

Chapter~\ref{chap:current_state} characterized the baseline operational context and documented current observability limitations, including fragmented telemetry and reactive troubleshooting patterns.

Chapter~\ref{chap:literature} synthesized the state of the art through systematic literature review following \ac{PRISMA} guidelines. The review identified relevant observability patterns, instrumentation techniques, and evaluation approaches documented in prior work. The synthesis established that while observability frameworks for cloud-native environments are extensively documented, evidence and guidance for operationally constrained, heterogeneous on-premise \ac{CICD} settings---often involving privilege-limited patterns even when partial Docker-level access exists---remain comparatively sparse, which motivates the design-to-evaluation work planned in subsequent chapters.

Chapter~\ref{chap:planning} defined methodology and requirements planning for execution. Chapter~\ref{chap:design} translated those requirements into tool selection criteria and architectural decisions. Chapter~\ref{chap:implementation} established the implementation strategy and quality gates for reproducible deployment. Chapter~\ref{chap:evaluation} specified the evaluation framework, including a-priori metrics and validity considerations. Together, these chapters form a coherent execution path from requirements to measurable outcomes.

\section{Transition to Thesis Execution}

The thesis execution phase will implement and evaluate the planned solution. Immediate next steps include:

\begin{itemize}
	\item Complete technology selection based on compatibility with identified constraints, documented in assessment matrix with explicit rationale.
	\item Design reference architecture specifying components, deployment topology, data flows, and correlation mechanisms.
	\item Implement prototype in staging environment, producing deployment templates, configuration artifacts, and operational procedures.
	\item Execute evaluation campaign comparing baseline and instrumented deployments under defined workloads with controlled fault injection.
	\item Collect and analyze measurement data, prepare dissertation chapters documenting implementation, evaluation findings, and implications.
\end{itemize}

Validation targets defined in project planning include quantitative metrics (\ac{MTTD}, \ac{MTTR}), telemetry pipeline throughput verification, alert precision and detection coverage assessment via fault injection, and instrumentation overhead quantification with an operational target below 5\% where feasible. These targets are operationalized through specific measurement procedures and acceptance criteria documented in the evaluation protocol. Analysis is restricted to tested contexts without overgeneralization.

The work documented in this \ac{PREPD} confirms readiness to proceed with thesis execution. All preparatory activities required to initiate implementation and evaluation have been completed.

